\documentclass[aps,preprint]{revtex4}%
\usepackage{amsfonts}
\usepackage{amsmath}
\usepackage{amssymb}
\usepackage{graphicx}%
\setcounter{MaxMatrixCols}{30}
%TCIDATA{OutputFilter=latex2.dll}
%TCIDATA{Version=4.10.0.2363}
%TCIDATA{CSTFile=revtex4.cst}
%TCIDATA{Created=Sunday, April 20, 2003 13:13:14}
%TCIDATA{LastRevised=Monday, May 19, 2003 11:16:43}
%TCIDATA{<META NAME="GraphicsSave" CONTENT="32">}
%TCIDATA{<META NAME="DocumentShell" CONTENT="Articles\SW\REVTeX 4">}
%TCIDATA{Language=American English}
\newtheorem{theorem}{Theorem}
\newtheorem{acknowledgement}[theorem]{Acknowledgement}
\newtheorem{algorithm}[theorem]{Algorithm}
\newtheorem{axiom}[theorem]{Axiom}
\newtheorem{claim}[theorem]{Claim}
\newtheorem{conclusion}[theorem]{Conclusion}
\newtheorem{condition}[theorem]{Condition}
\newtheorem{conjecture}[theorem]{Conjecture}
\newtheorem{corollary}[theorem]{Corollary}
\newtheorem{criterion}[theorem]{Criterion}
\newtheorem{definition}[theorem]{Definition}
\newtheorem{example}[theorem]{Example}
\newtheorem{exercise}[theorem]{Exercise}
\newtheorem{lemma}[theorem]{Lemma}
\newtheorem{notation}[theorem]{Notation}
\newtheorem{problem}[theorem]{Problem}
\newtheorem{proposition}[theorem]{Proposition}
\newtheorem{remark}[theorem]{Remark}
\newtheorem{solution}[theorem]{Solution}
\newtheorem{summary}[theorem]{Summary}
\newenvironment{proof}[1][Proof]{\noindent\textbf{#1.} }{\ \rule{0.5em}{0.5em}}
\begin{document}
\preprint{ }
\title[Adapted Operator Basis]{Contraction Adapted Operator Basis}
\author{Brian Weiner}
\affiliation{PSU}
\keywords{}
\pacs{}

\begin{abstract}


\end{abstract}
\date[Date text]{date}
\tableofcontents


\section{Operator Space Basis}

\subsection{Generalized Bases for $\mathcal{B}\left(  \mathcal{H}^{N}\right)
$ and $\mathcal{B}_{1}\left(  \mathcal{H}^{N}\right)  $}

The duality between $\mathcal{B}_{1}\left(  \mathcal{H}^{N}\right)  $ and
$\mathcal{B}\left(  \mathcal{H}^{N}\right)  $ can be expressed in terms of the
scalar product
\begin{equation}
\left(  X|Y\right)  =Tr\left\{  X^{\dagger}Y\right\}  ;\quad X\in
\mathcal{B}\left(  \mathcal{H}^{N}\right)  ,Y\in\mathcal{B}_{1}\left(
\mathcal{H}^{N}\right)
\end{equation}
Following the Dirac convention we can denote operators, $Y\in\mathcal{B}%
_{1}\left(  \mathcal{H}^{N}\right)  $, as $\left\vert Y\right)  $ and
operators, $X\in\mathcal{B}\left(  \mathcal{H}^{N}\right)  $, as $\left(
X\right\vert $ thus producing the super operator%
\begin{equation}
\left\vert Y\right)  \left(  X\right\vert :\mathcal{B}_{1}\left(
\mathcal{H}^{N}\right)  \rightarrow\mathcal{B}_{1}\left(  \mathcal{H}%
^{N}\right)
\end{equation}
whose action is given by
\begin{equation}
\left\vert Y\right)  \left(  X\right\vert \,Z=Tr\left\{  X^{\dagger}Z\right\}
\left\vert Y\right)  \equiv Tr\left\{  X^{\dagger}Z\right\}  Y
\end{equation}


We can introduce a generalized basis, $\left\{  \left\vert \mathsf{z}%
^{N}\right\rangle \left\langle \mathsf{z}^{\prime N}\right\vert \right\}  $,
\emph{for both} of the dual spaces $\mathcal{B}\left(  \mathcal{H}^{N}\right)
$ and $\mathcal{B}_{1}\left(  \mathcal{H}^{N}\right)  $, composed of ket-bra
operators that correspond to the generalized eigenvectors of the $N$-particle
space-spin coordinate operator. None of the basis elements $\left\vert
\mathsf{z}^{N}\right\rangle \left\langle \mathsf{z}^{\prime N}\right\vert $
belong to $\mathcal{B}\left(  \mathcal{H}^{N}\right)  $ as the action of any
individual $\left\vert \mathsf{z}^{N}\right\rangle \left\langle \mathsf{z}%
^{\prime N}\right\vert $ on any vector belonging to $\mathcal{H}^{N}$ produces
a vector $\notin\mathcal{H}^{N}$. The notation $\left\vert \left\vert
\mathsf{z}^{N}\right\rangle \left\langle \mathsf{z}^{\prime N}\right\vert
\right)  $ signifies that $\left\vert \mathsf{z}^{N}\right\rangle \left\langle
\mathsf{z}^{\prime N}\right\vert $ is to be considered within the
$\mathcal{B}_{1}\left(  \mathcal{H}^{N}\right)  $-topology, while the notation
$\left(  \left\vert \mathsf{z}^{N}\right\rangle \left\langle \mathsf{z}%
^{\prime N}\right\vert \right\vert $ signifies that $\left\vert \mathsf{z}%
^{N}\right\rangle \left\langle \mathsf{z}^{\prime N}\right\vert $ is to be
considered within the $\mathcal{B}\left(  \mathcal{H}^{N}\right)  $-topology.

This basis allows one to define a super operator resolution of the identity on
$\mathcal{B}_{1}\left(  \mathcal{H}^{N}\right)  $ by
\begin{equation}
\hat{I}_{\mathcal{B}_{1}\left(  \mathcal{H}^{N}\right)  }=%
%TCIMACRO{\dint }%
%BeginExpansion
{\displaystyle\int}
%EndExpansion
\left\vert \left\vert \mathsf{z}^{N}\right\rangle \left\langle \mathsf{z}%
^{\prime N}\right\vert \right)  \left(  \left\vert \mathsf{z}^{N}\right\rangle
\left\langle \mathsf{z}^{\prime N}\right\vert \right\vert d\mathsf{z}%
^{N}d\mathsf{z}^{\prime N},
\end{equation}
which can be expressed symbolically as
\begin{equation}
\hat{I}_{\mathcal{B}_{1}\left(  \mathcal{H}^{N}\right)  }=\left(  \left(
\underline{\left\vert \mathsf{z}^{N}\right\rangle \left\langle \mathsf{z}%
^{\prime N}\right\vert }\right\vert \right)  \left(  \left\vert \underline
{\left\vert \mathsf{z}^{N}\right\rangle \left\langle \mathsf{z}^{\prime
N}\right\vert }\right)  \right)  ^{t},
\end{equation}
producing the generalized expansion of an element $X\in\mathcal{B}_{1}\left(
\mathcal{H}^{N}\right)  $ as%
\begin{align}
X  &  =%
%TCIMACRO{\dint }%
%BeginExpansion
{\displaystyle\int}
%EndExpansion
\left\vert \left\vert \mathsf{z}^{N}\right\rangle \left\langle \mathsf{z}%
^{\prime N}\right\vert \right)  \left(  \left\vert \mathsf{z}^{N}\right\rangle
\left\langle \mathsf{z}^{\prime N}\right\vert |X\right)  d\mathsf{z}%
^{N}d\mathsf{z}^{\prime N}\nonumber\\
&  =%
%TCIMACRO{\dint }%
%BeginExpansion
{\displaystyle\int}
%EndExpansion
\left\vert \left\vert \mathsf{z}^{N}\right\rangle \left\langle \mathsf{z}%
^{\prime N}\right\vert \right)  Tr\left\{  \left\vert \mathsf{z}%
^{N}\right\rangle \left\langle \mathsf{z}^{\prime N}\right\vert ^{\dagger
}X\right\}  d\mathsf{z}^{N}d\mathsf{z}^{\prime N}\nonumber\\
&  =%
%TCIMACRO{\dint }%
%BeginExpansion
{\displaystyle\int}
%EndExpansion
\left\vert \left\vert \mathsf{z}^{N}\right\rangle \left\langle \mathsf{z}%
^{\prime N}\right\vert \right)  \left\langle \mathsf{z}^{N}|X\mathsf{z}%
^{\prime N}\right\rangle d\mathsf{z}^{N}d\mathsf{z}^{\prime N}\nonumber\\
&  =%
%TCIMACRO{\dint }%
%BeginExpansion
{\displaystyle\int}
%EndExpansion
X\left(  \mathsf{z}^{N},\mathsf{z}^{\prime N}\right)  \left\vert
\mathsf{z}^{N}\right\rangle \left\langle \mathsf{z}^{\prime N}\right\vert
d\mathsf{z}^{N}d\mathsf{z}^{\prime N}%
\end{align}
where $X\left(  \mathsf{z}^{N},\mathsf{z}^{\prime N}\right)  $ is the integral
kernel corresponding to $X$.

The space of Hilbert-Schmidt operators $\mathcal{B}_{2}\left(  \mathcal{H}%
^{N}\right)  $ is a Hilbert space and is contained in the Banach space
$\mathcal{B}\left(  \mathcal{H}^{N}\right)  $. The predual $\mathcal{B}%
_{1}\left(  \mathcal{H}^{N}\right)  $ of $\mathcal{B}\left(  \mathcal{H}%
^{N}\right)  $ is contained in $\mathcal{B}_{2}\left(  \mathcal{H}^{N}\right)
,$ thus $\mathcal{B}\left(  \mathcal{H}^{N}\right)  \supset\mathcal{B}%
_{2}\left(  \mathcal{H}^{N}\right)  \supset\mathcal{B}_{1}\left(
\mathcal{H}^{N}\right)  $ form a Gelfand Triple all spanned by the basis
$\left\{  \left\vert \mathsf{z}^{N}\right\rangle \left\langle \mathsf{z}%
^{\prime N}\right\vert \right\}  $.

The basis $\left\{  \left\vert \mathsf{z}^{N}\right\rangle \left\langle
\mathsf{z}^{\prime N}\right\vert \right\}  $ can be graded into sub bases
$\left\{  \mathfrak{b}_{m}\left(  \mathsf{z}\right)  \right\}  $ with
precisely $m$ bra and ket coordinates \emph{unequal} which are defined as
\begin{equation}
\mathfrak{b}_{m}\left(  \mathsf{z}\right)  =\left\{  \left\vert \mathsf{z}%
^{N}\right\rangle \left\langle \mathsf{z}^{\prime N}\right\vert |\mathsf{z}%
^{N}\overset{m}{\neq}\mathsf{z}^{\prime N}\right\}
\end{equation}
where $\mathsf{z}^{N}\overset{m}{\neq}\mathsf{z}^{\prime N}$ denotes that the
coordinates $\mathsf{z}^{N}$ and $\mathsf{z}^{\prime N}$ have exactly $m$
values different, which could be equivalently written as $\mathsf{z}%
^{N}\overset{N-m}{=}\mathsf{z}^{\prime N}$ i.e. $\mathsf{z}^{N}$ and
$\mathsf{z}^{\prime N}$ have exactly $N-m$ values in common. These bases
produce a graded decomposition of $\mathcal{B}_{1}\left(  \mathcal{H}%
^{N}\right)  $ as the direct sum
\begin{equation}
\mathcal{B}_{1}\left(  \mathcal{H}^{N}\right)  =%
%TCIMACRO{\dbigoplus \limits_{0\leq m\leq N}}%
%BeginExpansion
{\displaystyle\bigoplus\limits_{0\leq m\leq N}}
%EndExpansion
\mathfrak{B}_{m}\left(  \mathsf{z}^{N}\right)
\end{equation}
where $\mathfrak{B}_{m}\left(  \mathsf{z}\right)  $ is defined as the closed
linear span of the basis $\left\{  \mathfrak{b}_{m}\left(  \mathsf{z}\right)
\right\}  $ i.e.
\begin{equation}
\mathfrak{B}_{m}\left(  \mathsf{z}^{N}\right)  =\overline
{\operatorname{linspan}}\left\{  \left.  \underset{}{\left\vert \mathsf{z}%
^{N}\right\rangle \left\langle \mathsf{z}^{\prime N}\right\vert }\right\vert
\left\vert \mathsf{z}^{N}\right\rangle \left\langle \mathsf{z}^{\prime
N}\right\vert \in\mathfrak{b}_{m}\left(  \mathsf{z}^{N}\right)  \right\}
\end{equation}
and the closure is taken wrt the trace norm topology. In order to analyze the
kernels of the contraction maps $\left\{  L_{N}^{p}|0\leq p\leq N\right\}  $
it is useful to consider the sequence of subspaces $\left\{  \mathfrak{V}%
_{p}\left(  \mathsf{z}\right)  \right\}  $, that are generated by $\left\{
\left\vert \mathsf{z}^{N}\right\rangle \left\langle \mathsf{z}^{\prime
N}\right\vert \right\}  $ that have at least $p$ differences between
$\mathsf{z}^{N}$ and $\mathsf{z}^{\prime N}$ (denoted by $\mathsf{z}%
^{N}\overset{p}{\nsim}$ $\mathsf{z}^{\prime N}$) or equivalently generated by
bases elements that have at most $N-p$ values in common between $\mathsf{z}%
^{N}$ and $\mathsf{z}^{\prime N}$ (denoted by $\mathsf{z}^{N}\overset
{N-p}{\sim}$ $\mathsf{z}^{\prime N})$, which are defined as
\begin{equation}
\mathfrak{V}_{p}\left(  \mathsf{z}^{N}\right)  =%
%TCIMACRO{\dbigoplus \limits_{p\leq m\leq N}}%
%BeginExpansion
{\displaystyle\bigoplus\limits_{p\leq m\leq N}}
%EndExpansion
\mathfrak{B}_{m}\left(  \mathsf{z}^{N}\right)
\end{equation}
so that
\begin{equation}
\mathfrak{V}_{p}\left(  \mathsf{z}^{N}\right)  =\overline
{\operatorname{linspan}}\left\{  \left.  \left\vert \mathsf{z}^{N}%
\right\rangle \left\langle \mathsf{z}^{\prime N}\right\vert \overset{}%
{}\right\vert \mathsf{z}^{N}\overset{p}{\nsim}\mathsf{z}^{\prime N}\right\}
=\overline{\operatorname{linspan}}\left\{  \mathfrak{c}_{p}\left(
\mathsf{z}^{N}\right)  \right\}
\end{equation}
where $\mathfrak{c}_{p}\left(  \mathsf{z}^{N}\right)  \equiv\left\{  \left.
\underset{}{\left\vert \mathsf{z}^{N}\right\rangle \left\langle \mathsf{z}%
^{\prime N}\right\vert }\right\vert \mathsf{z}^{N}\overset{p}{\nsim}%
\mathsf{z}^{\prime N}\right\}  $ $\equiv\left\{  \left.  \underset
{}{\left\vert \mathsf{z}^{N}\right\rangle \left\langle \mathsf{z}^{\prime
N}\right\vert }\right\vert \mathsf{z}^{N}\overset{N-p}{\sim}\mathsf{z}^{\prime
N}\right\}  $ and hence $\mathfrak{V}_{0}\left(  \mathsf{z}^{N}\right)
=\mathcal{B}_{1}\left(  \mathcal{H}^{N}\right)  $. One should also observe
that
\begin{equation}
\mathfrak{B}_{p}\left(  \mathsf{z}^{N}\right)  =\mathfrak{V}_{p-1}\left(
\mathsf{z}^{N}\right)  \ominus\mathfrak{V}_{p}\left(  \mathsf{z}^{N}\right)
\end{equation}
These subspaces partially characterize the kernel of the maps $\left\{
L_{N}^{p}\right\}  $ as they have the property
\begin{equation}
\mathfrak{V}_{p}\left(  \mathsf{z}^{N}\right)  \supset\ker L_{N}^{p}%
\supset\mathfrak{V}_{p+1}\left(  \mathsf{z}^{N}\right)  .
\end{equation}
and one has in general%
\begin{equation}
\mathcal{B}_{1}\left(  \mathcal{H}^{N}\right)  =\mathfrak{V}_{0}\left(
\mathsf{z}^{N}\right)  \supset\ker L_{N}^{0}\supset\mathfrak{V}_{1}\left(
\mathsf{z}\right)  \supset\cdots\supset\ker L_{N}^{N-1}\supset\mathfrak{V}%
_{N}\left(  \mathsf{z}^{N}\right)  \supset\ker L_{N}^{N}=\phi
\end{equation}


\subsection{Proper Bases for $\mathcal{B}\left(  \mathcal{H}^{N}\right)  $ and
$\mathcal{B}_{1}\left(  \mathcal{H}^{N}\right)  $}

One can repeat all the above for proper bases of $\mathcal{H}^{N}$ generated
by complete orthonormal bases $\left\{  \varphi_{j}\right\}  $ of
$\mathcal{H}^{N}$ to produce $\mathfrak{B}_{p}\left(  \mathbf{\varphi}%
^{N}\right)  ,$ $\mathfrak{b}_{p}\left(  \mathbf{\varphi}^{N}\right)
,\mathfrak{c}_{p}\left(  \mathbf{\varphi}^{N}\right)  $ and $\mathfrak{V}%
_{p}\left(  \mathbf{\varphi}^{N}\right)  $ and their associated decompositions
of $\mathcal{B}_{1}\left(  \mathcal{H}^{N}\right)  .$ Hence in particular
\begin{equation}
\mathfrak{V}_{p}\left(  \mathbf{\varphi}^{N}\right)  \supset\ker L_{N}%
^{p}\supset\mathfrak{V}_{p+1}\left(  \mathbf{\varphi}^{N}\right)
\end{equation}
The space with at \emph{least} $p$ differences (with \emph{most N-p} the
same)
\begin{equation}
\mathfrak{V}_{p}\left(  \mathbf{\varphi}^{N}\right)  =%
%TCIMACRO{\dbigoplus \limits_{p\leq m\leq N}}%
%BeginExpansion
{\displaystyle\bigoplus\limits_{p\leq m\leq N}}
%EndExpansion
\mathfrak{B}_{m}\left(  \mathbf{\varphi}^{N}\right)
\end{equation}
The space with $m$ different coordinates can be expressed as
\begin{equation}
\mathfrak{B}_{m}\left(  \mathbf{\varphi}^{N}\right)  =\overline
{\operatorname{linspan}}\left\{  \left.  \underset{}{\left\vert
\mathbf{\varphi}^{N}\right\rangle \left\langle \varphi^{\prime N}\right\vert
}\right\vert \left\vert \mathbf{\varphi}^{N}\right\rangle \left\langle
\varphi^{\prime N}\right\vert \in\mathfrak{b}_{m}\left(  \mathbf{\varphi}%
^{N}\right)  \right\}
\end{equation}
and the space with at least $p$ differences can be expressed as%
\begin{equation}
\mathfrak{V}_{p}\left(  \mathbf{\varphi}^{N}\right)  =\overline
{\operatorname{linspan}}\left\{  \left.  \left\vert \mathbf{\varphi}%
^{N}\right\rangle \left\langle \varphi^{\prime N}\right\vert \overset{}%
{}\right\vert \mathbf{\varphi}^{N}\overset{p}{\nsim}\varphi^{\prime
N}\right\}  =\overline{\operatorname{linspan}}\left\{  \mathfrak{c}_{p}\left(
\mathbf{\varphi}^{N}\right)  \right\}
\end{equation}
The space with exactly $p$ different coordinates can be expressed as the
difference of the spaces with at least $p-1$ differences and the space with at
least $p$ differences
\begin{equation}
\mathfrak{B}_{p}\left(  \mathbf{\varphi}^{N}\right)  =\mathfrak{V}%
_{p-1}\left(  \mathbf{\varphi}^{N}\right)  \ominus\mathfrak{V}_{p}\left(
\mathbf{\varphi}^{N}\right)
\end{equation}


\section{Representations of $U\left(  \mathcal{H}^{1}\right)  $}

One can define subspaces of $\mathcal{B}_{1}\left(  \mathcal{H}^{N}\right)  $
that produce irreducible representations of $SU\left(  \mathcal{H}^{1}\right)
$ that are based on the kernels of the contraction maps $\left\{  \left.
L_{N}^{P}\right\vert 0\leq p\leq N\right\}  ,$ namely
\begin{equation}
\mathcal{K}_{Np}=\ker L_{N}^{p-1}\ominus\ker L_{N}^{p};\quad1\leq p\leq N
\end{equation}
and
\begin{equation}
\mathcal{B}_{1}\left(  \mathcal{H}^{N}\right)  =%
%TCIMACRO{\tbigoplus \limits_{1\leq p\leq N}}%
%BeginExpansion
{\textstyle\bigoplus\limits_{1\leq p\leq N}}
%EndExpansion
\mathcal{K}_{Np}\oplus\mathcal{K}_{N0}%
\end{equation}
where $\mathcal{K}_{N0}$ is a compliment of $%
%TCIMACRO{\tbigoplus \limits_{1\leq p\leq N}}%
%BeginExpansion
{\textstyle\bigoplus\limits_{1\leq p\leq N}}
%EndExpansion
\mathcal{K}_{Np}$ in $\mathcal{B}_{1}\left(  \mathcal{H}^{N}\right)  $. These
subspaces are related to the ones introduced previously by $\mathcal{K}%
_{Np}\supset\mathfrak{B}_{p+1}\left(  \mathbf{\varphi}^{N}\right)  $.

\subsection{One Particle Operators $\mathcal{B}_{1}\left(  \mathcal{H}%
^{1}\right)  $}

In particular one can consider the bases of the operator space $\mathcal{B}%
_{1}\left(  \mathcal{H}^{1}\right)  $ given by%
\begin{equation}
\left\{  \left\{  \left.  \left\vert \varphi_{j}\right\rangle \left\langle
\varphi_{j}\right\vert -\left\vert \varphi_{j+1}\right\rangle \left\langle
\varphi_{j+1}\right\vert \right\vert 1\leq j\leq r\right\}  ,\left\{  \left.
\left\vert \varphi_{j}\right\rangle \left\langle \varphi_{k}\right\vert
\right\vert 1\leq j\neq k\leq r\right\}  ,\left\vert \varphi_{1}\right\rangle
\left\langle \varphi_{1}\right\vert \right\}
\end{equation}
that is adapted to the decomposition
\begin{equation}
\mathcal{B}_{1}\left(  \mathcal{H}^{1}\right)  =\mathcal{K}_{11od}%
\oplus\mathcal{K}_{11d}\oplus\mathcal{K}_{10}\label{graddec}%
\end{equation}
where
\begin{align}
\mathcal{K}_{11od} &  =\overline{\operatorname{linspan}}\left\{  \left.
\left\vert \varphi_{j}\right\rangle \left\langle \varphi_{k}\right\vert
\right\vert 1\leq j\neq k\leq r\right\}  \nonumber\\
\mathcal{K}_{11d} &  =\overline{\operatorname{linspan}}\left\{  \left.
\left\vert \varphi_{j}\right\rangle \left\langle \varphi_{j}\right\vert
-\left\vert \varphi_{j+1}\right\rangle \left\langle \varphi_{j+1}\right\vert
\right\vert 1\leq j\leq r\right\}  \nonumber\\
\mathcal{K}_{10} &  =\left\{  \left.  \alpha\left\vert \varphi_{1}%
\right\rangle \left\langle \varphi_{1}\right\vert \right\vert \alpha
\in\mathbb{C}\right\}
\end{align}
Let us denote these basis elements by
\begin{align}
\left\vert \varphi_{jk}\right)   &  =\left\vert \varphi_{j}\right\rangle
\left\langle \varphi_{k}\right\vert ;\quad1\leq j\neq k\leq r\nonumber\\
\left\vert \varphi_{j+1}\right)   &  =\left\vert \varphi_{j}\right\rangle
\left\langle \varphi_{j}\right\vert -\left\vert \varphi_{j+1}\right\rangle
\left\langle \varphi_{j+1}\right\vert ;\quad1\leq j\leq r\nonumber\\
\left\vert \varphi_{1}\right)   &  =\left\vert \varphi_{1}\right\rangle
\left\langle \varphi_{1}\right\vert
\end{align}
We can note that
\begin{equation}
\left\vert \varphi_{1}\right\rangle \left\langle \varphi_{1}\right\vert
-\left(  \left\vert \varphi_{1}\right\rangle \left\langle \varphi
_{1}\right\vert -\left\vert \varphi_{2}\right\rangle \left\langle \varphi
_{2}\right\vert \right)  =\left\vert \varphi_{2}\right\rangle \left\langle
\varphi_{2}\right\vert
\end{equation}
and
\begin{equation}
\left(  \left\vert \varphi_{1}\right\rangle \left\langle \varphi
_{1}\right\vert -\left\vert \varphi_{2}\right\rangle \left\langle \varphi
_{2}\right\vert \right)  +\left(  \left\vert \varphi_{2}\right\rangle
\left\langle \varphi_{2}\right\vert -\left\vert \varphi_{3}\right\rangle
\left\langle \varphi_{3}\right\vert \right)  =\left(  \left\vert \varphi
_{1}\right\rangle \left\langle \varphi_{1}\right\vert -\left\vert \varphi
_{3}\right\rangle \left\langle \varphi_{3}\right\vert \right)
\end{equation}
so that
\begin{equation}
\left\vert \varphi_{1}\right\rangle \left\langle \varphi_{1}\right\vert
-\left(  \left\vert \varphi_{1}\right\rangle \left\langle \varphi
_{1}\right\vert -\left\vert \varphi_{2}\right\rangle \left\langle \varphi
_{2}\right\vert \right)  -\left(  \left\vert \varphi_{2}\right\rangle
\left\langle \varphi_{2}\right\vert -\left\vert \varphi_{3}\right\rangle
\left\langle \varphi_{3}\right\vert \right)  =\left\vert \varphi
_{3}\right\rangle \left\langle \varphi_{3}\right\vert
\end{equation}
etc., hence generating all diagonal basis elements $\left\{  \left.
\left\vert \varphi_{j}\right\rangle \left\langle \varphi_{j}\right\vert
\right\vert ;2\leq j\leq r\right\}  $ as%
\begin{equation}
\left\vert \varphi_{j}\right\rangle \left\langle \varphi_{j}\right\vert
=\left\vert \varphi_{1}\right)  -\sum_{2\leq m\leq j}\left\vert \varphi
_{j}\right)  ;\quad1\leq j\leq r
\end{equation}
Any element of $X\in\mathcal{B}_{1}\left(  \mathcal{H}^{1}\right)  $ can be
expanded%
\begin{align}
A &  =\sum_{1\leq j,k\leq\infty}A_{jk}\left\vert \varphi_{j}\right\rangle
\left\langle \varphi_{k}\right\vert \\
&  =\sum_{1\leq j\neq k\leq r}X_{jk}\left\vert \varphi_{jk}\right)
+\sum_{2\leq j\leq r}X_{j}\left\vert \varphi_{j}\right)  +X_{1}\left\vert
\varphi_{1}\right)
\end{align}
where
\begin{align}
X_{1} &  =Tr\left\{  A\right\}  \nonumber\\
X_{jk} &  =A_{jk}\nonumber\\
X_{j} &  =Tr\left\{  A\right\}  -\sum_{1\leq m\leq j-1}A_{mm};\quad2\leq j\leq
r
\end{align}
The choice of $\left\vert \varphi_{1}\right)  =\left\vert \varphi
_{1}\right\rangle \left\langle \varphi_{1}\right\vert $ is quite arbitrary,
one can just as well use $\left\vert \varphi_{1}\right)  =\left\vert
\varphi_{j}\right\rangle \left\langle \varphi_{j}\right\vert $ for any other
$j,$ which just produces a change of basis for $\mathcal{B}_{1}\left(
\mathcal{H}^{N}\right)  $ and hence a change of expansion coefficients while
still produces a decomposition that is isomorphic to eq. $\left(
\ref{graddec}\right)  $.

\subsection{Two Particle Operators $\mathcal{B}_{1}\left(  \mathcal{H}%
^{2}\right)  $}

$\mathcal{B}_{1}\left(  \mathcal{H}^{2}\right)  $ can be decomposed as
\begin{align}
\mathcal{B}_{1}\left(  \mathcal{H}^{2}\right)   &  =\left(  \ker L_{2}%
^{1}\ominus\ker L_{2}^{2}\right)  \oplus\left(  \ker L_{2}^{0}\ominus\ker
L_{2}^{1}\right)  \oplus\left(  \ker L_{2}^{0}\right)  ^{c}\nonumber\\
&  =\ker L_{2}^{0}\oplus\left(  \ker L_{2}^{0}\right)  ^{c}\nonumber\\
&  =\ker L_{2}^{1}\oplus\left(  \ker L_{2}^{1}\right)  ^{c}\label{expan1}\\
&  =\mathcal{K}_{22}\oplus\mathcal{K}_{21}\oplus\mathcal{K}_{20}%
\end{align}
as $\ker L_{2}^{2}=\phi$. It can be shown that if $L_{2}^{1}$ is surjective
then
\begin{equation}
\mathcal{K}_{21}\oplus\mathcal{K}_{20}\cong\mathcal{B}_{1}\left(
\mathcal{H}^{1}\right)
\end{equation}
One can show that $L_{p}^{p-1}$ is surjective if $\dim\left(  \mathcal{H}%
^{1}\right)  >2(p-1).$ The decomposition eq. $\left(  \ref{expan1}\right)  $
can be analyzed in terms of the lie groups as
\begin{equation}
\mathcal{B}_{1}\left(  \mathcal{H}^{2}\right)  =sl\left(  \mathcal{H}%
^{2}\right)  _{2}\oplus sl\left(  \mathcal{H}^{2}\right)  _{2}^{c}=gl\left(
\mathcal{H}^{2}\right)  _{2}%
\end{equation}
as
\begin{align}
sl\left(  \mathcal{H}^{2}\right)  _{2} &  =\mathcal{K}_{22}\oplus
\mathcal{K}_{21}\nonumber\\
sl\left(  \mathcal{H}^{2}\right)  _{2}^{c} &  =\mathcal{K}_{20}%
\end{align}
and
\begin{equation}
\mathcal{B}_{1}\left(  \mathcal{H}^{2}\right)  =\left(  sl\left(
\mathcal{H}^{2}\right)  _{2}\ominus sl\left(  \mathcal{H}^{1}\right)
_{2}\right)  \oplus sl\left(  \mathcal{H}^{1}\right)  _{2}\oplus sl\left(
\mathcal{H}^{2}\right)  _{2}^{c}%
\end{equation}
where $sl\left(  \mathcal{H}^{2}\right)  _{2}$ and $sl\left(  \mathcal{H}%
^{1}\right)  _{2}$ are representations of the special linear lie algebras
$sl\left(  \mathcal{H}^{2}\right)  $ and $sl\left(  \mathcal{H}^{1}\right)  $
respectively in $\mathcal{H}^{2}.$ By exponentiating these subspaces one
obtains representations of the lie groups
\begin{align}
SL\left(  \mathcal{H}^{2}\right)  _{2} &  =\exp\left(  sl\left(
\mathcal{H}^{2}\right)  _{2}\right)  \nonumber\\
SL\left(  \mathcal{H}^{1}\right)  _{2} &  =\exp\left(  sl\left(
\mathcal{H}^{1}\right)  _{2}\right)
\end{align}
The decomposition eq. $\left(  \ref{expan1}\right)  $ also has the property
that it carries irreducible adjoint representation of $SU\left(
\mathcal{H}^{1}\right)  $ i.e.%
\begin{align}
\operatorname{Ad}\left(  SU\left(  \mathcal{H}^{1}\right)  _{2}\right)   &
:\mathcal{K}_{22}\rightarrow\mathcal{K}_{22}\nonumber\\
\operatorname{Ad}\left(  SU\left(  \mathcal{H}^{1}\right)  _{2}\right)   &
:\mathcal{K}_{21}\rightarrow\mathcal{K}_{21}%
\end{align}
which can be deduced from the intertwining property
\begin{equation}
L_{N}^{P}\circ\operatorname{Ad}\left(  SU\left(  \mathcal{H}^{1}\right)
_{p}\right)  =\operatorname{Ad}\left(  SU\left(  \mathcal{H}^{1}\right)
_{p}\right)  \circ L_{N}^{P}%
\end{equation}
which gives that
\begin{equation}
\operatorname{Ad}\left(  SU\left(  \mathcal{H}^{1}\right)  _{p}\right)  :\ker
L_{N}^{P}\rightarrow\ker L_{N}^{P}%
\end{equation}
so in particular
\begin{equation}
L_{2}^{1}\circ\operatorname{Ad}\left(  SU\left(  \mathcal{H}^{1}\right)
_{1}\right)  =\operatorname{Ad}\left(  SU\left(  \mathcal{H}^{1}\right)
_{1}\right)  \circ L_{2}^{1}%
\end{equation}
The subspaces $\mathcal{K}_{22}$ and $\mathcal{K}_{21}$ can be decomposed into
diagonal $\left\{  \left\vert \varphi_{j_{1}}\varphi_{j_{2}}\right\rangle
\left\langle \varphi_{j_{1}}\varphi_{j_{2}}\right\vert \right\}  $ and non
diagonal $\left\{  \left.  \left\vert \varphi_{j_{1}}\varphi_{j_{2}%
}\right\rangle \left\langle \varphi_{k_{1}}\varphi_{k_{2}}\right\vert
\right\vert j_{1},j_{2}\neq k_{1},k_{2}\right\}  $ parts wrt the basis
$\left\{  \left\vert \varphi_{j_{1}}\varphi_{j_{2}}\right\rangle \right\}  $
of $\mathcal{H}^{2}$ i.e.
\begin{align}
\mathcal{K}_{22} &  =\mathcal{K}_{22od}\oplus\mathcal{K}_{22d}\nonumber\\
\mathcal{K}_{21} &  =\mathcal{K}_{21od}\oplus\mathcal{K}_{21d}%
\end{align}
so that
\begin{align}
\mathcal{B}_{1}\left(  \mathcal{H}^{2}\right)   &  =\mathcal{K}_{22od}%
\oplus\mathcal{K}_{22d}\oplus\mathcal{K}_{21od}\oplus\mathcal{K}_{21d}%
\oplus\mathcal{K}_{20}\nonumber\\
&  =\ker L_{2}^{1}\oplus\left(  \ker L_{2}^{1}\right)  ^{c}\nonumber\\
&  =\left(  \ker L_{2}^{1}\ominus\ker L_{2}^{2}\right)  \oplus\left(  \ker
L_{2}^{0}\ominus\ker L_{2}^{1}\right)  \oplus\left(  \ker L_{2}^{0}\right)
^{c}\nonumber\\
&  =\mathcal{K}_{22}\oplus\mathcal{K}_{21}\oplus\mathcal{K}_{20}%
\end{align}
One can define diagonal and off-diagonal two particle operator subspaces as
\begin{align}
\mathcal{K}_{2d} &  =\mathcal{K}_{22d}\oplus\mathcal{K}_{21d}\nonumber\\
\mathcal{K}_{2od} &  =\mathcal{K}_{22od}\oplus\mathcal{K}_{21od}%
\end{align}
so that
\begin{equation}
\mathcal{B}_{1}\left(  \mathcal{H}^{2}\right)  =\mathcal{K}_{2d}%
\oplus\mathcal{K}_{2od}%
\end{equation}
and observe that
\begin{align}
L_{2}^{1} &  :\mathcal{K}_{2d}\rightarrow\mathcal{K}_{1d}\nonumber\\
L_{2}^{1} &  :\mathcal{K}_{2od}\rightarrow\mathcal{K}_{1od}%
\end{align}
which can seen by considering for example
\begin{align}
L_{2}^{1}\left(  \left\vert \varphi_{1}\varphi_{2}\right\rangle \left\langle
\varphi_{1}\varphi_{3}\right\vert \right)   &  =\left\vert \varphi
_{2}\right\rangle \left\langle \varphi_{3}\right\vert \nonumber\\
L_{2}^{1}\left(  \left\vert \varphi_{1}\varphi_{2}\right\rangle \left\langle
\varphi_{1}\varphi_{2}\right\vert \right)   &  =\left\vert \varphi
_{1}\right\rangle \left\langle \varphi_{1}\right\vert +\left\vert \varphi
_{2}\right\rangle \left\langle \varphi_{2}\right\vert
\end{align}


\subsubsection{Example $\dim\left(  \mathcal{H}^{1}\right)  =4$}

We can consider as an example a space $\mathcal{H}^{1}$ that has dimension 4.
In this case the number of different coordinates of the basis of
$\mathcal{B}_{1}\left(  \mathcal{H}^{1}\right)  $ can be displayed by the
matrix
\begin{equation}
\left(
\begin{array}
[c]{ccccccc}%
\ast & 12 & 13 & 14 & 23 & 24 & 34\\
12 & 0 & 1 & 1 & 1 & 1 & 2\\
13 & 1 & 0 & 1 & 1 & 2 & 1\\
14 & 1 & 1 & 0 & 2 & 1 & 1\\
23 & 1 & 1 & 2 & 0 & 1 & 1\\
24 & 1 & 2 & 1 & 1 & 0 & 1\\
34 & 2 & 1 & 1 & 1 & 1 & 0
\end{array}
\right)
\end{equation}


Then
\begin{align}
&  \dim\left(  \mathcal{H}^{1}\right)  =r=4;\quad\dim\left(  \mathcal{H}%
^{2}\right)  =\binom{r}{2}=\binom{4}{2}=6\nonumber\\
&  \dim\left(  \mathcal{B}_{1}\left(  \mathcal{H}^{1}\right)  \right)
=r^{2}=16;\quad\dim\left(  \mathcal{B}_{1}\left(  \mathcal{H}^{2}\right)
\right)  =\binom{r}{2}^{2}=36\nonumber\\
&  \dim\left(  \mathcal{K}_{21}\oplus\mathcal{K}_{20}\right)  =\dim\left(
\mathcal{B}_{1}\left(  \mathcal{H}^{1}\right)  \right)  =16\nonumber\\
&  \dim\left(  \mathcal{K}_{2d}\right)  =\binom{r}{2}=6\nonumber\\
&  \dim\left(  \mathcal{K}_{21d}\right)  =r-1=3\nonumber\\
&  \dim\left(  \mathcal{K}_{21od}\right)  =r\left(  r-1\right)  =12\nonumber\\
&  \dim\left(  \mathcal{K}_{20}\right)  =1\nonumber\\
&  \dim\left(  \mathcal{K}_{2}\right)  =\binom{r}{2}\left(  \binom{r}%
{2}-1\right)  =30\nonumber\\
&  \dim\left(  \mathcal{K}_{22}\right)  =\dim\left(  \mathcal{B}_{1}\left(
\mathcal{H}^{2}\right)  \right)  -\dim\left(  \mathcal{B}_{1}\left(
\mathcal{H}^{1}\right)  \right)  =36-16=20\nonumber\\
&  \dim\left(  \mathcal{K}_{22od}\right)  =\binom{r}{2}\binom{r-2}{2}%
+2\binom{r}{2}\left(  \frac{r-2}{2}\right)  +2\binom{r}{2}\left(  \frac
{r-4}{2}\right) \nonumber\\
&  =\binom{4}{2}\binom{2}{2}+2\binom{2}{2}\left(  \frac{2}{2}\right)
+2\binom{2}{2}\left(  \frac{0}{2}\right)  =6\times1+2\times6\times
1=18\nonumber\\
&  \dim\left(  \mathcal{K}_{22d}\right)  =\dim\left(  \mathcal{K}_{2d}\right)
-\dim\left(  \mathcal{K}_{21d}\right)  =6-4=2
\end{align}


The space $\mathfrak{B}_{2}\left(  \mathbf{\varphi}^{N}\right)  \subset
\mathcal{B}_{1}\left(  \mathcal{H}^{2}\right)  $ with no coordinates in common
has dimension
\begin{equation}
\binom{r}{2}\binom{r-2}{2}=\binom{4}{2}\binom{2}{2}=\left(  6\right)  \left(
1\right)  =6
\end{equation}
these basis elements have the form
\begin{equation}
\left\vert \varphi_{j_{1}}\right\rangle \left\langle \varphi_{k_{1}%
}\right\vert \wedge\left\vert \varphi_{j_{2}}\right\rangle \left\langle
\varphi_{k_{2}}\right\vert ;\quad1\leq\left(  j_{1}<j_{2}\right)  \neq\left(
k_{1}<k_{2}\right)  \leq4
\end{equation}
The space $\mathfrak{B}_{1}\left(  \mathbf{\varphi}^{N}\right)  \subset
\mathcal{B}_{1}\left(  \mathcal{H}^{2}\right)  $ with one difference has
dimension
\begin{equation}
2\binom{r}{2}\left(  r-2\right)  =2\left(  6\right)  \left(  2\right)  =24
\end{equation}
these basis elements have the form
\begin{equation}
\left\vert \varphi_{j_{1}}\right\rangle \left\langle \varphi_{k_{1}%
}\right\vert \wedge\left\vert \varphi_{j_{2}}\right\rangle \left\langle
\varphi_{j_{2}}\right\vert ;\quad1\leq j_{1}\neq k_{1}\leq4;\quad1\leq
j_{2}\neq j_{1},k_{1}\leq4
\end{equation}
The space $\mathfrak{B}_{0}\left(  \mathbf{\varphi}^{N}\right)  \subset
\mathcal{B}_{1}\left(  \mathcal{H}^{2}\right)  $ with no differences has
dimension%
\begin{equation}
\binom{r}{2}=\binom{4}{2}=6
\end{equation}
these basis elements have the form
\begin{equation}
\left\vert \varphi_{j_{1}}\right\rangle \left\langle \varphi_{j_{1}%
}\right\vert \wedge\left\vert \varphi_{j_{2}}\right\rangle \left\langle
\varphi_{j_{2}}\right\vert ;\quad1\leq j_{1}\neq j_{2}\leq4
\end{equation}
The sum of the dimensions is
\begin{equation}
6+24+6=36
\end{equation}
which is the dimension of $\mathcal{B}_{1}\left(  \mathcal{H}^{2}\right)  .$

The space $\mathcal{K}_{22od}$ is generated by the $6$ linear independent
elements which include the basis of $\mathfrak{B}_{2}\left(  \mathbf{\varphi
}^{2}\right)  $
\begin{equation}
\mathfrak{b}_{2}\left(  \mathbf{\varphi}^{2}\right)  =\left\{  \left\vert
\varphi_{1}\varphi_{2}\right\rangle \left\langle \varphi_{3}\varphi
_{4}\right\vert ,\left\vert \varphi_{1}\varphi_{3}\right\rangle \left\langle
\varphi_{2}\varphi_{4}\right\vert ,\left\vert \varphi_{1}\varphi
_{4}\right\rangle \left\langle \varphi_{2}\varphi_{3}\right\vert ,\left\vert
\varphi_{2}\varphi_{3}\right\rangle \left\langle \varphi_{1}\varphi
_{4}\right\vert ,\left\vert \varphi_{2}\varphi_{4}\right\rangle \left\langle
\varphi_{1}\varphi_{3}\right\vert ,\left\vert \varphi_{3}\varphi
_{4}\right\rangle \left\langle \varphi_{1}\varphi_{2}\right\vert \right\}
\end{equation}
$\lceil$In general
\begin{align}
&  \left\{  \left\vert \varphi_{j_{1}k_{1}}\right)  \wedge\left\vert
\varphi_{j_{2}k_{2}}\right)  =\left\vert \varphi_{j_{1}}\right\rangle
\left\langle \varphi_{k_{1}}\right\vert \wedge\left\vert \varphi_{j_{2}%
}\right\rangle \left\langle \varphi_{k_{2}}\right\vert =\left\vert
\varphi_{j_{1}}\varphi_{j_{2}}\right\rangle \left\langle \varphi_{k_{1}%
}\varphi_{k_{2}}\right\vert \right. \nonumber\\
&  \qquad\qquad\qquad\qquad\left.  j_{1}\neq k_{1},k_{2};\quad j_{2}\neq
k_{1},k_{2};\quad j_{1}<j_{2};k_{1}<k_{2}\right\}
\end{align}
and contains $\binom{r}{2}\times\binom{r-2}{2}$ basis elements$\rfloor$ and
the subspace of $\mathfrak{B}_{1}\left(  \mathbf{\varphi}^{2}\right)  $
generated by the basis
\begin{align}
&  \left\{  \left(  \left\vert \varphi_{1}\right\rangle \left\langle
\varphi_{1}\right\vert -\left\vert \varphi_{2}\right\rangle \left\langle
\varphi_{2}\right\vert \right)  \wedge\left\vert \varphi_{3}\right\rangle
\left\langle \varphi_{4}\right\vert ,\left(  \left\vert \varphi_{1}%
\right\rangle \left\langle \varphi_{1}\right\vert -\left\vert \varphi
_{2}\right\rangle \left\langle \varphi_{2}\right\vert \right)  \wedge
\left\vert \varphi_{4}\right\rangle \left\langle \varphi_{3}\right\vert
\right. \nonumber\\
&  \,\,\,\left(  \left\vert \varphi_{1}\right\rangle \left\langle \varphi
_{1}\right\vert -\left\vert \varphi_{3}\right\rangle \left\langle \varphi
_{3}\right\vert \right)  \wedge\left\vert \varphi_{2}\right\rangle
\left\langle \varphi_{4}\right\vert ,\left(  \left\vert \varphi_{1}%
\right\rangle \left\langle \varphi_{1}\right\vert -\left\vert \varphi
_{3}\right\rangle \left\langle \varphi_{3}\right\vert \right)  \wedge
\left\vert \varphi_{2}\right\rangle \left\langle \varphi_{4}\right\vert
\nonumber\\
&  \,\,\,\left(  \left\vert \varphi_{1}\right\rangle \left\langle \varphi
_{1}\right\vert -\left\vert \varphi_{4}\right\rangle \left\langle \varphi
_{4}\right\vert \right)  \wedge\left\vert \varphi_{2}\right\rangle
\left\langle \varphi_{3}\right\vert ,\left(  \left\vert \varphi_{1}%
\right\rangle \left\langle \varphi_{1}\right\vert -\left\vert \varphi
_{4}\right\rangle \left\langle \varphi_{4}\right\vert \right)  \wedge
\left\vert \varphi_{3}\right\rangle \left\langle \varphi_{2}\right\vert
\nonumber\\
&  \,\,\,\left(  \left\vert \varphi_{2}\right\rangle \left\langle \varphi
_{2}\right\vert -\left\vert \varphi_{3}\right\rangle \left\langle \varphi
_{3}\right\vert \right)  \wedge\left\vert \varphi_{1}\right\rangle
\left\langle \varphi_{4}\right\vert ,\left(  \left\vert \varphi_{2}%
\right\rangle \left\langle \varphi_{2}\right\vert -\left\vert \varphi
_{3}\right\rangle \left\langle \varphi_{3}\right\vert \right)  \wedge
\left\vert \varphi_{4}\right\rangle \left\langle \varphi_{1}\right\vert
\nonumber\\
&  \,\,\,\left(  \left\vert \varphi_{2}\right\rangle \left\langle \varphi
_{2}\right\vert -\left\vert \varphi_{4}\right\rangle \left\langle \varphi
_{4}\right\vert \right)  \wedge\left\vert \varphi_{1}\right\rangle
\left\langle \varphi_{3}\right\vert ,\left(  \left\vert \varphi_{2}%
\right\rangle \left\langle \varphi_{2}\right\vert -\left\vert \varphi
_{4}\right\rangle \left\langle \varphi_{4}\right\vert \right)  \wedge
\left\vert \varphi_{3}\right\rangle \left\langle \varphi_{1}\right\vert
\nonumber\\
&  \,\left.  \left(  \left\vert \varphi_{3}\right\rangle \left\langle
\varphi_{3}\right\vert -\left\vert \varphi_{4}\right\rangle \left\langle
\varphi_{4}\right\vert \right)  \wedge\left\vert \varphi_{1}\right\rangle
\left\langle \varphi_{2}\right\vert ,\left(  \left\vert \varphi_{3}%
\right\rangle \left\langle \varphi_{3}\right\vert -\left\vert \varphi
_{4}\right\rangle \left\langle \varphi_{4}\right\vert \right)  \wedge
\left\vert \varphi_{2}\right\rangle \left\langle \varphi_{1}\right\vert
\right\}  \subset\mathfrak{b}_{1}\left(  \mathbf{\varphi}^{2}\right)
\end{align}
$\lceil$In general
\begin{equation}
\left(  \left\vert \varphi_{j_{1}}\right\rangle \left\langle \varphi_{j_{1}%
}\right\vert -\left\vert \varphi_{j_{2}}\right\rangle \left\langle
\varphi_{j_{2}}\right\vert \right)  \wedge\left\vert \varphi_{j_{3}%
}\right\rangle \left\langle \varphi_{j_{4}}\right\vert ;1\leq j_{1}<j_{2}%
\neq\left(  j_{3}\neq j_{4}\right)
\end{equation}
and contains $2\binom{r}{2}\binom{r-2}{2}$ elements$\rfloor.$

The space $\mathcal{K}_{22d}$ is generated by a subset of diagonal basis
elements
\begin{align}
&  \left\{  \left(  \left\vert \varphi_{1}\right\rangle \left\langle
\varphi_{1}\right\vert -\left\vert \varphi_{2}\right\rangle \left\langle
\varphi_{2}\right\vert \right)  \wedge\left(  \left\vert \varphi
_{3}\right\rangle \left\langle \varphi_{3}\right\vert -\left\vert \varphi
_{4}\right\rangle \left\langle \varphi_{4}\right\vert \right)  ,\right.
\nonumber\\
&  \left.  \left(  \left\vert \varphi_{1}\right\rangle \left\langle
\varphi_{1}\right\vert -\left\vert \varphi_{3}\right\rangle \left\langle
\varphi_{3}\right\vert \right)  \wedge\left(  \left\vert \varphi
_{2}\right\rangle \left\langle \varphi_{2}\right\vert -\left\vert \varphi
_{4}\right\rangle \left\langle \varphi_{4}\right\vert \right)  \right\}
\subset\mathfrak{b}_{0}\left(  \mathbf{\varphi}^{2}\right)
\end{align}
$\mathbf{\lceil}$In general
\begin{align}
&  \text{ }\left\{  \left(  \left\vert \varphi_{j_{1}}\right\rangle
\left\langle \varphi_{j_{1}}\right\vert -\left\vert \varphi_{j_{2}%
}\right\rangle \left\langle \varphi_{j_{2}}\right\vert \right)  \wedge\left(
\left\vert \varphi_{j_{3}}\right\rangle \left\langle \varphi_{j_{3}%
}\right\vert -\left\vert \varphi_{j_{4}}\right\rangle \left\langle
\varphi_{j_{4}}\right\vert \right)  \right.  =\nonumber\\
&  \,\left\vert \varphi_{j_{1}}\varphi_{j_{3}}\right\rangle \left\langle
\varphi_{j_{1}}\varphi_{j_{3}}\right\vert -\left\vert \varphi_{j_{1}}%
\varphi_{j_{4}}\right\rangle \left\langle \varphi_{j_{1}}\varphi_{j_{4}%
}\right\vert -\left\vert \varphi_{j_{2}}\varphi_{j_{3}}\right\rangle
\left\langle \varphi_{j_{2}}\varphi_{j_{3}}\right\vert +\left\vert
\varphi_{j_{2}}\varphi_{j_{4}}\right\rangle \left\langle \varphi_{j_{2}%
}\varphi_{j_{4}}\right\vert \quad\left.  j_{1}\neq j_{2}\neq j_{3}\neq
j_{4}\right\}
\end{align}
and contains $\binom{r}{2}^{2}-\binom{r}{2}\times\binom{r-2}{2}-r\left(
r-1\right)  $ elements$\rfloor$ to generate $\mathcal{K}_{22d}$ which is
evidently of dimension $2.$

We can check the result of the contraction map $L_{2}^{1}$ on these elements
by evaluating generic sample elements
\begin{align}
Tr\left\{  \left(  \left\vert \varphi_{1}\varphi_{3}\right\rangle \left\langle
\varphi_{1}\varphi_{4}\right\vert -\left\vert \varphi_{2}\varphi
_{3}\right\rangle \left\langle \varphi_{2}\varphi_{4}\right\vert \right)
a_{j}^{\dagger}a_{k}\right\}   &  =0\quad\forall j,k\nonumber\\
Tr\left\{  \left\vert \varphi_{1}\varphi_{2}\right\rangle \left\langle
\varphi_{3}\varphi_{4}\right\vert a_{j}^{\dagger}a_{k}\right\}   &
=0\quad\forall j,k\nonumber\\
Tr\left\{  \left(  \left\vert \varphi_{1}\varphi_{3}\right\rangle \left\langle
\varphi_{1}\varphi_{3}\right\vert -\left\vert \varphi_{1}\varphi
_{4}\right\rangle \left\langle \varphi_{1}\varphi_{4}\right\vert -\left\vert
\varphi_{2}\varphi_{3}\right\rangle \left\langle \varphi_{2}\varphi
_{3}\right\vert +\left\vert \varphi_{2}\varphi_{4}\right\rangle \left\langle
\varphi_{2}\varphi_{4}\right\vert \right)  a_{j}^{\dagger}a_{k}\right\}   &
=0\quad\forall j,k
\end{align}
Obviously such elements $\notin\ker L_{2}^{2}$ $=\phi$ as $\left\{  Tr\left\{
Ba_{j_{1}}^{\dagger}a_{j_{2}}^{\dagger}a_{k_{2}}a_{k_{1}}\right\}  \right\}  $
are not all zero.

The space $\mathcal{K}_{21od}$ of dimension $12$ is a subspace of
$\mathfrak{B}_{1}\left(  \mathbf{\varphi}^{2}\right)  $ and is generated by
\begin{align}
&  \left\{  \left(  \left\vert \varphi_{1}\right\rangle \left\langle
\varphi_{1}\right\vert +\left\vert \varphi_{2}\right\rangle \left\langle
\varphi_{2}\right\vert \right)  \wedge\left\vert \varphi_{3}\right\rangle
\left\langle \varphi_{4}\right\vert ,\left(  \left\vert \varphi_{1}%
\right\rangle \left\langle \varphi_{1}\right\vert +\left\vert \varphi
_{2}\right\rangle \left\langle \varphi_{2}\right\vert \right)  \wedge
\left\vert \varphi_{4}\right\rangle \left\langle \varphi_{3}\right\vert
\right. \nonumber\\
&  \left(  \,\,\,\left\vert \varphi_{1}\right\rangle \left\langle \varphi
_{1}\right\vert +\left\vert \varphi_{3}\right\rangle \left\langle \varphi
_{3}\right\vert \right)  \wedge\left\vert \varphi_{2}\right\rangle
\left\langle \varphi_{4}\right\vert ,\left(  \left\vert \varphi_{1}%
\right\rangle \left\langle \varphi_{1}\right\vert +\left\vert \varphi
_{3}\right\rangle \left\langle \varphi_{3}\right\vert \right)  \wedge
\left\vert \varphi_{2}\right\rangle \left\langle \varphi_{4}\right\vert
\nonumber\\
&  \left(  \,\,\,\left\vert \varphi_{1}\right\rangle \left\langle \varphi
_{1}\right\vert +\left\vert \varphi_{4}\right\rangle \left\langle \varphi
_{4}\right\vert \right)  \wedge\left\vert \varphi_{2}\right\rangle
\left\langle \varphi_{3}\right\vert ,\left(  \left\vert \varphi_{1}%
\right\rangle \left\langle \varphi_{1}\right\vert +\left\vert \varphi
_{4}\right\rangle \left\langle \varphi_{4}\right\vert \right)  \wedge
\left\vert \varphi_{3}\right\rangle \left\langle \varphi_{2}\right\vert
\nonumber\\
&  \left(  \,\,\,\left\vert \varphi_{2}\right\rangle \left\langle \varphi
_{2}\right\vert +\left\vert \varphi_{3}\right\rangle \left\langle \varphi
_{3}\right\vert \right)  \wedge\left\vert \varphi_{1}\right\rangle
\left\langle \varphi_{4}\right\vert ,\left(  \left\vert \varphi_{2}%
\right\rangle \left\langle \varphi_{2}\right\vert +\left\vert \varphi
_{3}\right\rangle \left\langle \varphi_{3}\right\vert \right)  \wedge
\left\vert \varphi_{4}\right\rangle \left\langle \varphi_{1}\right\vert
\nonumber\\
&  \left(  \,\,\,\left\vert \varphi_{2}\right\rangle \left\langle \varphi
_{2}\right\vert +\left\vert \varphi_{4}\right\rangle \left\langle \varphi
_{4}\right\vert \right)  \wedge\left\vert \varphi_{1}\right\rangle
\left\langle \varphi_{3}\right\vert ,\left(  \left\vert \varphi_{2}%
\right\rangle \left\langle \varphi_{2}\right\vert +\left\vert \varphi
_{4}\right\rangle \left\langle \varphi_{4}\right\vert \right)  \wedge
\left\vert \varphi_{3}\right\rangle \left\langle \varphi_{1}\right\vert
\nonumber\\
&  \left.  \left(  \,\left\vert \varphi_{3}\right\rangle \left\langle
\varphi_{3}\right\vert +\left\vert \varphi_{4}\right\rangle \left\langle
\varphi_{4}\right\vert \right)  \wedge\left\vert \varphi_{1}\right\rangle
\left\langle \varphi_{2}\right\vert ,\left(  \left\vert \varphi_{3}%
\right\rangle \left\langle \varphi_{3}\right\vert +\left\vert \varphi
_{4}\right\rangle \left\langle \varphi_{4}\right\vert \right)  \wedge
\left\vert \varphi_{2}\right\rangle \left\langle \varphi_{1}\right\vert
\right\}  \subset\mathfrak{b}_{1}\left(  \mathbf{\varphi}^{2}\right)
\end{align}
$\mathbf{\lceil}$In general
\begin{align}
&  \left\{  \left(  \left\vert \varphi_{j_{1}}\right\rangle \left\langle
\varphi_{j_{1}}\right\vert +\left\vert \varphi_{j_{2}}\right\rangle
\left\langle \varphi_{j_{2}}\right\vert \right)  \wedge\left\vert
\varphi_{j_{3}}\right\rangle \left\langle \varphi_{j_{4}}\right\vert \right.
=\nonumber\\
&  \,\left\vert \varphi_{j_{1}}\varphi_{j_{3}}\right\rangle \left\langle
\varphi_{j_{1}}\varphi_{j_{4}}\right\vert +\left\vert \varphi_{j_{2}}%
\varphi_{j_{3}}\right\rangle \left\langle \varphi_{j_{2}}\varphi_{j_{4}%
}\right\vert \qquad\left.  \left(  j_{1}<j_{2}\right)  \neq\left(  j_{3}\neq
j_{4}\right)  \right\}
\end{align}
and contains $r\left(  r-1\right)  $ elements$\rfloor.$

The space $\mathcal{K}_{21d}$ must be of dimension $r-1=3$ and has a basis
given by
\begin{align}
&  \left\{  \left\vert \varphi_{1}\right\rangle \left\langle \varphi
_{1}\right\vert \wedge\left(  \left\vert \varphi_{2}\right\rangle \left\langle
\varphi_{2}\right\vert -\left\vert \varphi_{3}\right\rangle \left\langle
\varphi_{3}\right\vert \right)  ,\right. \nonumber\\
&  \,\,\nonumber\\
&  \left.  \,\left\vert \varphi_{1}\right\rangle \left\langle \varphi
_{1}\right\vert \wedge\left(  \left\vert \varphi_{3}\right\rangle \left\langle
\varphi_{3}\right\vert -\left\vert \varphi_{4}\right\rangle \left\langle
\varphi_{4}\right\vert \right)  \right\}  \subset\mathfrak{b}_{0}\left(
\mathbf{\varphi}^{2}\right)
\end{align}
$\mathbf{\lceil}$In general
\begin{equation}
\left\{  \left\vert \varphi_{1}\right\rangle \left\langle \varphi
_{1}\right\vert \wedge\left(  \left\vert \varphi_{j_{1}}\right\rangle
\left\langle \varphi_{j_{1}}\right\vert -\left\vert \varphi_{j_{2}%
}\right\rangle \left\langle \varphi_{j_{2}}\right\vert \right)  \right.
=\left\vert \varphi_{1}\varphi_{j_{1}}\right\rangle \left\langle \varphi
_{1}\varphi_{j_{1}}\right\vert -\left\vert \varphi_{1}\varphi_{j_{2}%
}\right\rangle \left\langle \varphi_{1}\varphi_{j_{2}}\right\vert
\qquad\left.  j_{1}\neq j_{2}\neq1\right\}
\end{equation}
and contains $r-1$ elements$\rfloor.$

and $\dim\left(  \mathcal{K}_{20}\right)  =1$ and is spanned by
\begin{equation}
\left\vert \varphi_{1}\varphi_{2}\right\rangle \left\langle \varphi_{1}%
\varphi_{2}\right\vert
\end{equation}


We can check the result of the contraction map $L_{2}^{0}$ on these elements
by evaluating
\begin{equation}
Tr\left\{  B\right\}
\end{equation}
which gives for example
\begin{align*}
Tr\left\{  \left\vert \varphi_{1}\varphi_{2}\right\rangle \left\langle
\varphi_{1}\varphi_{3}\right\vert \right\}   &  =0\\
Tr\left\{  \left\vert \varphi_{1}\varphi_{2}\right\rangle \left\langle
\varphi_{1}\varphi_{2}\right\vert -\left\vert \varphi_{1}\varphi
_{3}\right\rangle \left\langle \varphi_{1}\varphi_{3}\right\vert \right\}   &
=0
\end{align*}
and the results of contraction map $L_{2}^{1}$ by
\begin{equation}
Tr\left\{  Ba_{j}^{\dagger}a_{k}\right\}
\end{equation}
which gives for example%
\begin{align*}
Tr\left\{  \left\vert \varphi_{1}\varphi_{2}\right\rangle \left\langle
\varphi_{1}\varphi_{3}\right\vert a_{j}^{\dagger}a_{k}\right\}   &
=0\quad\not \forall \,j,k\\
Tr\left\{  \left(  \left\vert \varphi_{1}\varphi_{2}\right\rangle \left\langle
\varphi_{1}\varphi_{2}\right\vert -\left\vert \varphi_{1}\varphi
_{3}\right\rangle \left\langle \varphi_{1}\varphi_{3}\right\vert \right)
a_{j}^{\dagger}a_{k}\right\}   &  =0\quad\not \forall \,j,k
\end{align*}


\subsubsection{Example of $\dim\left(  \mathcal{H}^{1}\right)  =6$}

We can consider as an example a space $\mathcal{H}^{1}$ that has dimension
$6$. In this case the number of different coordinates of the basis of
$\mathcal{B}_{1}\left(  \mathcal{H}^{1}\right)  $ can be displayed by the
matrix
\begin{equation}
\left(
\begin{array}
[c]{cccccccccccccccc}%
\ast & 12 & 13 & 14 & 15 & 16 & 23 & 24 & 25 & 26 & 34 & 35 & 36 & 45 & 46 &
56\\
12 & 0 & 1 & 1 & 1 & 1 & 1 & 1 & 1 & 1 & 2 & 2 & 2 & 2 & 2 & 2\\
13 & 1 & 0 & 1 & 1 & 1 & 1 & 2 & 2 & 2 & 1 & 1 & 1 & 2 & 2 & 2\\
14 & 1 & 1 & 0 & 1 & 1 & 2 & 1 & 2 & 2 & 1 & 2 & 2 & 1 & 1 & 2\\
15 & 1 & 1 & 1 & 0 & 1 & 2 & 2 & 1 & 2 & 2 & 1 & 2 & 1 & 2 & 1\\
16 & 1 & 1 & 1 & 1 & 0 & 2 & 2 & 2 & 1 & 2 & 2 & 1 & 2 & 1 & 1\\
23 & 1 & 1 & 2 & 2 & 2 & 0 & 1 & 1 & 1 & 1 & 1 & 1 & 2 & 2 & 2\\
24 & 1 & 2 & 1 & 2 & 2 & 1 & 0 & 1 & 1 & 1 & 2 & 2 & 1 & 1 & 2\\
25 & 1 & 2 & 2 & 1 & 2 & 1 & 1 & 0 & 1 & 2 & 1 & 2 & 1 & 2 & 1\\
26 & 1 & 2 & 2 & 2 & 1 & 1 & 1 & 1 & 0 & 2 & 2 & 1 & 2 & 1 & 1\\
34 & 2 & 1 & 1 & 2 & 2 & 1 & 1 & 2 & 2 & 0 & 1 & 1 & 1 & 1 & 2\\
35 & 2 & 1 & 2 & 1 & 2 & 1 & 2 & 1 & 2 & 1 & 0 & 1 & 1 & 2 & 1\\
36 & 2 & 1 & 2 & 2 & 1 & 1 & 2 & 2 & 1 & 1 & 1 & 0 & 2 & 1 & 1\\
45 & 2 & 2 & 1 & 1 & 2 & 2 & 1 & 1 & 2 & 1 & 1 & 2 & 0 & 1 & 1\\
46 & 2 & 2 & 1 & 2 & 1 & 2 & 1 & 2 & 1 & 1 & 2 & 1 & 1 & 0 & 1\\
56 & 2 & 2 & 2 & 1 & 1 & 2 & 2 & 1 & 1 & 2 & 1 & 1 & 1 & 1 & 0
\end{array}
\right)
\end{equation}
Then
\begin{align}
&  \dim\left(  \mathcal{H}^{2}\right)  =r=6;\quad\dim\left(  \mathcal{H}%
^{2}\right)  =\binom{r}{2}=\binom{6}{2}=15\nonumber\\
&  \dim\left(  \mathcal{B}_{1}\left(  \mathcal{H}^{1}\right)  \right)
=r^{2}=36;\quad\dim\left(  \mathcal{B}_{1}\left(  \mathcal{H}^{2}\right)
\right)  =\binom{r}{2}^{2}=225\nonumber\\
&  \dim\left(  \mathcal{K}_{21}\oplus\mathcal{K}_{20}\right)  =\dim\left(
\mathcal{B}_{1}\left(  \mathcal{H}^{1}\right)  \right)  =36\nonumber\\
&  \dim\left(  \mathcal{K}_{21d}\right)  =5\nonumber\\
&  \dim\left(  \mathcal{K}_{21od}\right)  =r\left(  r-1\right)  =30\nonumber\\
&  \dim\left(  \mathcal{K}_{20}\right)  =1\nonumber\\
&  \dim\left(  \mathcal{K}_{22}\right)  =\dim\left(  \mathcal{B}_{1}\left(
\mathcal{H}^{2}\right)  \right)  -\dim\left(  \mathcal{B}_{1}\left(
\mathcal{H}^{1}\right)  \right)  =225-36=189\nonumber\\
&  \dim\left(  \mathcal{K}_{22od}\right)  =\binom{r}{2}\binom{r-2}{2}%
+2\binom{r}{2}\left(  \frac{r-2}{2}\right)  +2\binom{r}{2}\left(  \frac
{r-4}{2}\right) \nonumber\\
&  =\binom{6}{2}\binom{4}{2}+2\binom{6}{2}\left(  \frac{4}{2}\right)
+2\binom{6}{2}\left(  \frac{6-4}{2}\right)  =15\times6+2\times15\times
2+2\times15\times1=180\nonumber\\
&  \dim\left(  \mathcal{K}_{22d}\right)  =\binom{r}{2}-\binom{r}%
{1}=15-6=9\nonumber\\
&  \dim\left(  \mathcal{K}_{2d}\right)  =\binom{r}{2}=15\nonumber\\
&  \dim\left(  \mathcal{K}_{2od}\right)  =\binom{r}{2}\left(  \binom{r}%
{2}-1\right)  =210
\end{align}
The space $\mathfrak{B}_{2}\left(  \mathbf{\varphi}^{N}\right)  \subset
\mathcal{B}_{1}\left(  \mathcal{H}^{2}\right)  $ with no coordinates in common
has dimension
\begin{equation}
\binom{r}{2}\binom{r-2}{2}=\binom{6}{2}\binom{4}{2}=\left(  15\right)  \left(
6\right)  =90
\end{equation}
these basis elements have the form
\begin{equation}
\left\vert \varphi_{j_{1}}\right\rangle \left\langle \varphi_{k_{1}%
}\right\vert \wedge\left\vert \varphi_{j_{2}}\right\rangle \left\langle
\varphi_{k_{2}}\right\vert ;\quad1\leq\left(  j_{1}<j_{2}\right)  \neq\left(
k_{1}<k_{2}\right)  \leq6
\end{equation}
The space $\mathfrak{B}_{1}\left(  \mathbf{\varphi}^{N}\right)  \subset
\mathcal{B}_{1}\left(  \mathcal{H}^{2}\right)  $ with one difference has
dimension
\begin{equation}
2\binom{r}{2}\left(  r-2\right)  =2\left(  15\right)  \left(  4\right)  =120
\end{equation}
these basis elements have the form
\begin{equation}
\left\vert \varphi_{j_{1}}\right\rangle \left\langle \varphi_{k_{1}%
}\right\vert \wedge\left\vert \varphi_{j_{2}}\right\rangle \left\langle
\varphi_{j_{2}}\right\vert ;\quad1\leq j_{1}\neq k_{1}\leq6;\quad1\leq
j_{2}\neq j_{1},k_{1}\leq6
\end{equation}
The space $\mathfrak{B}_{0}\left(  \mathbf{\varphi}^{N}\right)  \subset
\mathcal{B}_{1}\left(  \mathcal{H}^{2}\right)  $ with no differences has
dimension%
\begin{equation}
\binom{r}{2}=\binom{6}{2}=15
\end{equation}
these basis elements have the form
\begin{equation}
\left\vert \varphi_{j_{1}}\right\rangle \left\langle \varphi_{j_{1}%
}\right\vert \wedge\left\vert \varphi_{j_{2}}\right\rangle \left\langle
\varphi_{j_{2}}\right\vert ;\quad1\leq j_{1}\neq j_{2}\leq6
\end{equation}
The sum of these dimensions is
\begin{equation}
90+120+15=225
\end{equation}
which is the dimension of $\mathcal{B}_{1}\left(  \mathcal{H}^{2}\right)  .$

The space $\mathfrak{B}_{2}\left(  \mathbf{\varphi}^{N}\right)  \subset
\mathcal{K}_{22}=\ker L_{2}^{1}$, while only parts of the spaces
$\mathfrak{B}_{1}\left(  \mathbf{\varphi}^{N}\right)  $ and $\mathfrak{B}%
_{0}\left(  \mathbf{\varphi}^{N}\right)  $ are contained in $\mathcal{K}_{22}%
$. The space $\mathfrak{B}_{2}\left(  \mathbf{\varphi}^{N}\right)
=\mathcal{K}_{22od}$ is generated by the $90$ linear independent elements
(i.e. is of dimension $90)$
\begin{equation}
\mathfrak{b}_{2}\left(  \mathbf{\varphi}^{N}\right)  =\left\vert \varphi
_{1}\varphi_{2}\right\rangle \left\langle \varphi_{3}\varphi_{4}\right\vert
,\cdots\cdots\cdots,\left\vert \varphi_{3}\varphi_{4}\right\rangle
\left\langle \varphi_{5}\varphi_{6}\right\vert
\end{equation}
$\lceil$In general
\begin{align}
&  \left\{  \left\vert \varphi_{j_{1}}\right\rangle \left\langle
\varphi_{k_{1}}\right\vert \wedge\left\vert \varphi_{j_{2}}\right\rangle
\left\langle \varphi_{k_{2}}\right\vert =\left\vert \varphi_{j_{1}}%
\varphi_{j_{2}}\right\rangle \left\langle \varphi_{k_{1}}\varphi_{k_{2}%
}\right\vert \right. \nonumber\\
&  \qquad\qquad\qquad\qquad\qquad\qquad\qquad\qquad\qquad\qquad\left.
j_{1}\neq k_{1},k_{2};\quad j_{2}\neq k_{1},k_{2};\quad j_{1}<j_{2}%
;k_{1}<k_{2}\right\}
\end{align}
and contains $\binom{r}{2}\binom{r-2}{2}$ basis elements$\rfloor$

The part of $\mathfrak{B}_{1}\left(  \mathbf{\varphi}^{N}\right)  $ that is
contained in $\mathcal{K}_{22od}\subset\mathcal{K}_{22}$ is formed by the
basis elements
\begin{equation}
\left\vert \varphi_{j_{1}}\right\rangle \left\langle \varphi_{k_{1}%
}\right\vert \wedge\left(  \left\vert \varphi_{j_{2}}\right\rangle
\left\langle \varphi_{j_{2}}\right\vert -\left\vert \varphi_{j_{3}%
}\right\rangle \left\langle \varphi_{j_{3}}\right\vert \right)  ;\quad
j_{1}\neq j_{2},j_{3};k_{1}\neq j_{2},j_{3};j_{2}<j_{3}%
\end{equation}
which number
\begin{equation}
2\binom{r}{2}\left(  \frac{r-2}{2}\right)  =2\left(  15\right)  \left(
2\right)  =60
\end{equation}
The corresponding plus combinations
\begin{equation}
\left\vert \varphi_{j_{1}}\right\rangle \left\langle \varphi_{k_{1}%
}\right\vert \wedge\left(  \left\vert \varphi_{j_{2}}\right\rangle
\left\langle \varphi_{j_{2}}\right\vert +\left\vert \varphi_{j_{3}%
}\right\rangle \left\langle \varphi_{j_{3}}\right\vert \right)  ;\quad
j_{1}\neq k_{1};j_{1}\neq j_{2},j_{3};k_{1}\neq j_{2},j_{3};j_{2}<j_{3}%
\end{equation}
number $60$, and some of these are also in $\mathcal{K}_{22od}\subset\ker
L_{2}^{1}$ as can be seen by considering $\left\{  \left\vert \varphi_{j_{1}%
}\right\rangle \left\langle \varphi_{k_{1}}\right\vert \wedge\left(
\left\vert \varphi_{j_{2}}\right\rangle \left\langle \varphi_{j_{2}%
}\right\vert +\left\vert \varphi_{j_{3}}\right\rangle \left\langle
\varphi_{j_{3}}\right\vert \right)  \right\}  $ for a fixed $\left(
j_{1},k_{1}\right)  ,$ which number $\left(  \frac{r-2}{2}\right)  =2.$ As
both of these map to $\left\vert \varphi_{j_{1}}\right\rangle \left\langle
\varphi_{k_{1}}\right\vert $ we can form $\pm$ combinations
\begin{equation}
\left\vert \varphi_{j_{1}}\right\rangle \left\langle \varphi_{k_{1}%
}\right\vert \wedge\left(  \left\vert \varphi_{j_{2}}\right\rangle
\left\langle \varphi_{j_{2}}\right\vert +\left\vert \varphi_{j_{3}%
}\right\rangle \left\langle \varphi_{j_{3}}\right\vert \right)  \pm\left\vert
\varphi_{j_{1}}\right\rangle \left\langle \varphi_{k_{1}}\right\vert
\wedge\left(  \left\vert \varphi_{j_{4}}\right\rangle \left\langle
\varphi_{j_{4}}\right\vert +\left\vert \varphi_{j_{5}}\right\rangle
\left\langle \varphi_{j_{5}}\right\vert \right)
\end{equation}
with one combination being in $\ker L_{2}^{1}$ and one being in its compliment
$\left(  \ker L_{2}^{1}\right)  ^{c}.$ The combinations%
\begin{align}
&  \left\vert \varphi_{j_{1}}\right\rangle \left\langle \varphi_{k_{1}%
}\right\vert \wedge\left(  \left\vert \varphi_{j_{2}}\right\rangle
\left\langle \varphi_{j_{2}}\right\vert +\left\vert \varphi_{j_{3}%
}\right\rangle \left\langle \varphi_{j_{3}}\right\vert \right)  -\left\vert
\varphi_{j_{1}}\right\rangle \left\langle \varphi_{k_{1}}\right\vert
\wedge\left(  \left\vert \varphi_{j_{4}}\right\rangle \left\langle
\varphi_{j_{4}}\right\vert +\left\vert \varphi_{j_{5}}\right\rangle
\left\langle \varphi_{j_{5}}\right\vert \right) \nonumber\\
&  j_{1}\neq k_{1};j_{1}\neq j_{2},j_{3};k_{1}\neq j_{2},j_{3};j_{2}%
<j_{3};j_{1}\neq k_{1};j_{1}\neq j_{4},j_{5};k_{1}\neq j_{4},j_{5};j_{4}<j_{5}%
\end{align}
are in $\mathcal{K}_{22od}\subset\ker L_{2}^{1}$ and number%
\begin{equation}
\binom{r}{2}\left(  \frac{r-2}{2}\right)  =30
\end{equation}
hence $\dim\mathcal{K}_{22od}=90+60+30=180.$ The plus combinations are in
$\left(  \ker L_{2}^{1}\right)  ^{c}$ and produce $\mathcal{K}_{1od}$ which is
of dimension $30$. The dimension of the space with a fixed $\left(
j_{1},k_{1}\right)  $ generated by
\begin{equation}
\left\vert \varphi_{j_{1}}\right\rangle \left\langle \varphi_{k_{1}%
}\right\vert \wedge\left(  \left\vert \varphi_{j_{2}}\right\rangle
\left\langle \varphi_{j_{2}}\right\vert +\left\vert \varphi_{j_{3}%
}\right\rangle \left\langle \varphi_{j_{3}}\right\vert \right)  ;\quad
j_{1}\neq k_{1};j_{1}\neq j_{2},j_{3};k_{1}\neq j_{2},j_{3};j_{2}<j_{3}%
\end{equation}
belonging to $\ker L_{2}^{1}$ is $\left(  \frac{r-2}{2}\right)  -1=\left(
\frac{r-4}{2}\right)  $ and the dimension of the space generated by all such
elements is given by
\begin{equation}
2\binom{r}{2}\left(  \frac{r-4}{2}\right)  =2\left(  15\right)  \left(
1\right)  =30
\end{equation}


The space $\mathcal{K}_{22d}$ is generated by elements of the type
\begin{align}
&  \left(  \left\vert \varphi_{1}\right\rangle \left\langle \varphi
_{1}\right\vert -\left\vert \varphi_{2}\right\rangle \left\langle \varphi
_{2}\right\vert \right)  \wedge\left(  \left\vert \varphi_{3}\right\rangle
\left\langle \varphi_{3}\right\vert -\left\vert \varphi_{4}\right\rangle
\left\langle \varphi_{4}\right\vert \right) \nonumber\\
&  \qquad\qquad\qquad=\left\vert \varphi_{1}\varphi_{3}\right\rangle
\left\langle \varphi_{1}\varphi_{3}\right\vert -\left\vert \varphi_{1}%
\varphi_{4}\right\rangle \left\langle \varphi_{1}\varphi_{4}\right\vert
-\left\vert \varphi_{2}\varphi_{3}\right\rangle \left\langle \varphi
_{2}\varphi_{3}\right\vert +\left\vert \varphi_{2}\varphi_{4}\right\rangle
\left\langle \varphi_{2}\varphi_{4}\right\vert \nonumber\\
&  \left(  \left\vert \varphi_{1}\right\rangle \left\langle \varphi
_{1}\right\vert -\left\vert \varphi_{2}\right\rangle \left\langle \varphi
_{2}\right\vert \right)  \wedge\left(  \left\vert \varphi_{5}\right\rangle
\left\langle \varphi_{5}\right\vert -\left\vert \varphi_{6}\right\rangle
\left\langle \varphi_{6}\right\vert \right) \nonumber\\
&  \qquad\qquad\qquad=\left\vert \varphi_{1}\varphi_{5}\right\rangle
\left\langle \varphi_{1}\varphi_{5}\right\vert -\left\vert \varphi_{1}%
\varphi_{6}\right\rangle \left\langle \varphi_{1}\varphi_{6}\right\vert
-\left\vert \varphi_{2}\varphi_{5}\right\rangle \left\langle \varphi
_{2}\varphi_{5}\right\vert +\left\vert \varphi_{2}\varphi_{6}\right\rangle
\left\langle \varphi_{2}\varphi_{6}\right\vert \nonumber\\
&  \left(  \left\vert \varphi_{3}\right\rangle \left\langle \varphi
_{3}\right\vert -\left\vert \varphi_{4}\right\rangle \left\langle \varphi
_{4}\right\vert \right)  \wedge\left(  \left\vert \varphi_{5}\right\rangle
\left\langle \varphi_{5}\right\vert -\left\vert \varphi_{6}\right\rangle
\left\langle \varphi_{6}\right\vert \right) \nonumber\\
&  \qquad\qquad\qquad=\left\vert \varphi_{3}\varphi_{5}\right\rangle
\left\langle \varphi_{3}\varphi_{5}\right\vert -\left\vert \varphi_{3}%
\varphi_{6}\right\rangle \left\langle \varphi_{3}\varphi_{6}\right\vert
-\left\vert \varphi_{4}\varphi_{5}\right\rangle \left\langle \varphi
_{4}\varphi_{5}\right\vert +\left\vert \varphi_{4}\varphi_{6}\right\rangle
\left\langle \varphi_{4}\varphi_{6}\right\vert \nonumber\\
&  \left(  \left\vert \varphi_{2}\right\rangle \left\langle \varphi
_{2}\right\vert -\left\vert \varphi_{3}\right\rangle \left\langle \varphi
_{3}\right\vert \right)  \wedge\left(  \left\vert \varphi_{4}\right\rangle
\left\langle \varphi_{4}\right\vert -\left\vert \varphi_{5}\right\rangle
\left\langle \varphi_{5}\right\vert \right) \nonumber\\
&  \qquad\qquad\qquad=\left\vert \varphi_{3}\varphi_{5}\right\rangle
\left\langle \varphi_{3}\varphi_{5}\right\vert -\left\vert \varphi_{3}%
\varphi_{4}\right\rangle \left\langle \varphi_{3}\varphi_{4}\right\vert
-\left\vert \varphi_{4}\varphi_{5}\right\rangle \left\langle \varphi
_{4}\varphi_{5}\right\vert +\left\vert \varphi_{2}\varphi_{4}\right\rangle
\left\langle \varphi_{2}\varphi_{4}\right\vert \nonumber\\
&  \left(  \left\vert \varphi_{2}\right\rangle \left\langle \varphi
_{2}\right\vert -\left\vert \varphi_{3}\right\rangle \left\langle \varphi
_{3}\right\vert \right)  \wedge\left(  \left\vert \varphi_{1}\right\rangle
\left\langle \varphi_{1}\right\vert -\left\vert \varphi_{6}\right\rangle
\left\langle \varphi_{6}\right\vert \right) \nonumber\\
&  \qquad\qquad\qquad=\left\vert \varphi_{1}\varphi_{2}\right\rangle
\left\langle \varphi_{1}\varphi_{2}\right\vert -\left\vert \varphi_{1}%
\varphi_{3}\right\rangle \left\langle \varphi_{1}\varphi_{3}\right\vert
-\left\vert \varphi_{2}\varphi_{6}\right\rangle \left\langle \varphi
_{2}\varphi_{6}\right\vert +\left\vert \varphi_{3}\varphi_{6}\right\rangle
\left\langle \varphi_{3}\varphi_{6}\right\vert \nonumber\\
&  \left(  \left\vert \varphi_{1}\right\rangle \left\langle \varphi
_{1}\right\vert -\left\vert \varphi_{6}\right\rangle \left\langle \varphi
_{6}\right\vert \right)  \wedge\left(  \left\vert \varphi_{4}\right\rangle
\left\langle \varphi_{4}\right\vert -\left\vert \varphi_{5}\right\rangle
\left\langle \varphi_{5}\right\vert \right) \nonumber\\
&  \qquad\qquad\qquad=\left\vert \varphi_{1}\varphi_{4}\right\rangle
\left\langle \varphi_{1}\varphi_{4}\right\vert -\left\vert \varphi_{1}%
\varphi_{5}\right\rangle \left\langle \varphi_{1}\varphi_{5}\right\vert
-\left\vert \varphi_{4}\varphi_{6}\right\rangle \left\langle \varphi
_{4}\varphi_{6}\right\vert +\left\vert \varphi_{5}\varphi_{6}\right\rangle
\left\langle \varphi_{5}\varphi_{6}\right\vert \nonumber\\
&  \left(  \left\vert \varphi_{1}\right\rangle \left\langle \varphi
_{1}\right\vert -\left\vert \varphi_{3}\right\rangle \left\langle \varphi
_{3}\right\vert \right)  \wedge\left(  \left\vert \varphi_{2}\right\rangle
\left\langle \varphi_{2}\right\vert -\left\vert \varphi_{4}\right\rangle
\left\langle \varphi_{4}\right\vert \right) \nonumber\\
&  \qquad\qquad\qquad=\left\vert \varphi_{1}\varphi_{2}\right\rangle
\left\langle \varphi_{1}\varphi_{2}\right\vert -\left\vert \varphi_{1}%
\varphi_{4}\right\rangle \left\langle \varphi_{1}\varphi_{4}\right\vert
-\left\vert \varphi_{2}\varphi_{3}\right\rangle \left\langle \varphi
_{2}\varphi_{3}\right\vert +\left\vert \varphi_{3}\varphi_{4}\right\rangle
\left\langle \varphi_{3}\varphi_{4}\right\vert \nonumber\\
&  \left(  \left\vert \varphi_{1}\right\rangle \left\langle \varphi
_{1}\right\vert -\left\vert \varphi_{3}\right\rangle \left\langle \varphi
_{3}\right\vert \right)  \wedge\left(  \left\vert \varphi_{4}\right\rangle
\left\langle \varphi_{4}\right\vert -\left\vert \varphi_{6}\right\rangle
\left\langle \varphi_{6}\right\vert \right) \nonumber\\
&  \qquad\qquad\qquad=\left\vert \varphi_{1}\varphi_{4}\right\rangle
\left\langle \varphi_{1}\varphi_{4}\right\vert -\left\vert \varphi_{1}%
\varphi_{6}\right\rangle \left\langle \varphi_{1}\varphi_{6}\right\vert
-\left\vert \varphi_{3}\varphi_{4}\right\rangle \left\langle \varphi
_{3}\varphi_{4}\right\vert +\left\vert \varphi_{3}\varphi_{6}\right\rangle
\left\langle \varphi_{3}\varphi_{6}\right\vert \nonumber\\
&  \left(  \left\vert \varphi_{2}\right\rangle \left\langle \varphi
_{2}\right\vert -\left\vert \varphi_{4}\right\rangle \left\langle \varphi
_{4}\right\vert \right)  \wedge\left(  \left\vert \varphi_{3}\right\rangle
\left\langle \varphi_{3}\right\vert -\left\vert \varphi_{5}\right\rangle
\left\langle \varphi_{5}\right\vert \right) \nonumber\\
&  \qquad\qquad\qquad=\left\vert \varphi_{2}\varphi_{3}\right\rangle
\left\langle \varphi_{2}\varphi_{3}\right\vert -\left\vert \varphi_{2}%
\varphi_{5}\right\rangle \left\langle \varphi_{2}\varphi_{5}\right\vert
-\left\vert \varphi_{3}\varphi_{4}\right\rangle \left\langle \varphi
_{3}\varphi_{4}\right\vert +\left\vert \varphi_{4}\varphi_{5}\right\rangle
\left\langle \varphi_{4}\varphi_{5}\right\vert
\end{align}
contains $9$

The space $\mathcal{K}_{21od}$ of dimension $30$ is generated by
\begin{align}
&  \left\{  \left(  \left\vert \varphi_{1}\right\rangle \left\langle
\varphi_{1}\right\vert +\left\vert \varphi_{2}\right\rangle \left\langle
\varphi_{2}\right\vert \right)  \wedge\left\vert \varphi_{3}\right\rangle
\left\langle \varphi_{4}\right\vert ,\left(  \left\vert \varphi_{1}%
\right\rangle \left\langle \varphi_{1}\right\vert +\left\vert \varphi
_{2}\right\rangle \left\langle \varphi_{2}\right\vert \right)  \wedge
\left\vert \varphi_{4}\right\rangle \left\langle \varphi_{3}\right\vert
\right. \nonumber\\
&  \left.  \left(  \,\left\vert \varphi_{5}\right\rangle \left\langle
\varphi_{5}\right\vert +\left\vert \varphi_{6}\right\rangle \left\langle
\varphi_{6}\right\vert \right)  \wedge\left\vert \varphi_{1}\right\rangle
\left\langle \varphi_{2}\right\vert ,\left(  \left\vert \varphi_{5}%
\right\rangle \left\langle \varphi_{5}\right\vert +\left\vert \varphi
_{6}\right\rangle \left\langle \varphi_{6}\right\vert \right)  \wedge
\left\vert \varphi_{2}\right\rangle \left\langle \varphi_{1}\right\vert
\right\}  \subset\mathfrak{b}_{1}\left(  \mathbf{\varphi}^{N}\right)
\end{align}
$\mathbf{\lceil}$In general
\begin{align}
&  \left\{  \left(  \left\vert \varphi_{j_{1}}\right\rangle \left\langle
\varphi_{j_{1}}\right\vert +\left\vert \varphi_{j_{2}}\right\rangle
\left\langle \varphi_{j_{2}}\right\vert \right)  \wedge\left\vert
\varphi_{j_{3}}\right\rangle \left\langle \varphi_{j_{4}}\right\vert \right.
=\nonumber\\
&  \,\left\vert \varphi_{j_{1}}\varphi_{j_{3}}\right\rangle \left\langle
\varphi_{j_{1}}\varphi_{j_{4}}\right\vert +\left\vert \varphi_{j_{2}}%
\varphi_{j_{3}}\right\rangle \left\langle \varphi_{j_{2}}\varphi_{j_{4}%
}\right\vert \qquad\left.  \left(  j_{1}<j_{2}\right)  \neq\left(  j_{3}\neq
j_{4}\right)  \right\}
\end{align}
and contains $r\left(  r-1\right)  $ elements$\rfloor.$

The space $\mathcal{K}_{21d}$ of dimension $5$ is generated by%
\begin{align}
&  \left\{  \left\vert \varphi_{1}\right\rangle \left\langle \varphi
_{1}\right\vert \wedge\left(  \left\vert \varphi_{2}\right\rangle \left\langle
\varphi_{2}\right\vert -\left\vert \varphi_{3}\right\rangle \left\langle
\varphi_{3}\right\vert \right)  ,\right. \nonumber\\
&  \,\,\,\left\vert \varphi_{1}\right\rangle \left\langle \varphi
_{1}\right\vert \wedge\left(  \left\vert \varphi_{3}\right\rangle \left\langle
\varphi_{3}\right\vert -\left\vert \varphi_{4}\right\rangle \left\langle
\varphi_{4}\right\vert \right) \nonumber\\
&  \,\left\vert \varphi_{1}\right\rangle \left\langle \varphi_{1}\right\vert
\wedge\left(  \left\vert \varphi_{4}\right\rangle \left\langle \varphi
_{4}\right\vert -\left\vert \varphi_{5}\right\rangle \left\langle \varphi
_{5}\right\vert \right) \nonumber\\
& \nonumber\\
&  \left.  \,\left\vert \varphi_{1}\right\rangle \left\langle \varphi
_{1}\right\vert \wedge\left(  \left\vert \varphi_{5}\right\rangle \left\langle
\varphi_{5}\right\vert -\left\vert \varphi_{6}\right\rangle \left\langle
\varphi_{6}\right\vert \right)  \right\}  \subset\mathfrak{b}_{0}\left(
\mathbf{\varphi}^{N}\right)
\end{align}


The space $\mathcal{K}_{20}$ of dimension $1$ is spanned by
\begin{equation}
\left\vert \varphi_{1}\varphi_{2}\right\rangle \left\langle \varphi_{1}%
\varphi_{2}\right\vert =\left\vert \varphi_{1}\right\rangle \left\langle
\varphi_{1}\right\vert \wedge\left\vert \varphi_{2}\right\rangle \left\langle
\varphi_{2}\right\vert \in\mathfrak{b}_{2}\left(  \mathbf{\varphi}^{N}\right)
\end{equation}



\end{document}